\gccFrontHeading{Abstract}

This thesis designs and plans a campus network solution for teaching, research, administration, and daily service scenarios in a university environment. Based on data communication technologies, the proposed solution adopts a hierarchical architecture consisting of core, aggregation, and access layers. It covers network topology design, IP address planning, routing policy, security control, and operation management. The design focuses on common campus-network requirements such as high-concurrency access, cross-area interconnection, resource access control, and future scalability, aiming to balance stability, security, and maintainability.

The study first analyzes the business requirements and constraints of campus network construction, including the network access characteristics of teaching areas, office areas, dormitory areas, and public service areas. It then proposes design methods based on mainstream network devices and protocols, including core switching, redundant links, virtual local area network segmentation, access control lists, and network address translation. Finally, the solution is evaluated from the perspectives of security, compatibility, scalability, and operating efficiency. The result indicates that the proposed design can support the basic services of a university campus network and provide a network foundation for future digital-campus construction.

\vspace{0.6em}
\noindent\textbf{Key words:} \gccKeywordsEn

